\documentclass[11pt,a4paper]{article}

\usepackage[utf8]{inputenc}
\usepackage[T1]{fontenc}
\usepackage[margin=2.5cm,top=3cm,bottom=3cm]{geometry}
\usepackage{xcolor}
\usepackage[unicode=true,colorlinks=true,linkcolor=blue!45!black,urlcolor=blue!45!black,
            pdftitle={Pesquisa Assistida por Inteligência Artificial — Dossiê de apresentação},
            pdfauthor={Euzébio Soares dos Santos}]{hyperref}

\setlength{\parindent}{0pt}
\setlength{\parskip}{5pt}
\renewcommand{\contentsname}{Sumário}

% ---------- macros simples ----------
\newcommand{\sub}[1]{{\bfseries #1}\par\smallskip}
\newcommand{\caixa}[1]{\par\smallskip\noindent\fbox{\parbox{0.97\textwidth}{\small #1}}\par\smallskip}
\newcommand{\fluxo}{\par\smallskip\noindent\texttt{CONCEITOS E HISTÓRIA $\rightarrow$ SETUP ADEQUADO $\rightarrow$ IA ADEQUADA $\rightarrow$ HIPÓTESE E MÉTODO}\par\smallskip}

\begin{document}

% ============================ CAPA ============================
\begin{titlepage}
\centering
\vspace*{2cm}
{\Huge\bfseries Pesquisa Assistida por\\[6pt] Inteligência Artificial}\\[14pt]
{\Large Ementa acadêmica e documento explicativo}\\[26pt]
{\large CONCEITOS \ \textbullet\ \ SETUP \ \textbullet\ \ IA ADEQUADA \ \textbullet\ \ INVESTIGAÇÃO}\\[34pt]
\begin{minipage}{0.86\textwidth}
\small
Uma disciplina prática para compreender um problema, preparar o ambiente computacional,
selecionar ferramentas e inteligências artificiais adequadas e construir uma investigação
verificável. Proposta-base para cursos de graduação, extensão e formação inicial em pesquisa.
\end{minipage}\\[34pt]
{\large\bfseries Dossiê de apresentação a instituição de ensino}\\[18pt]
{\large Euzébio Soares dos Santos}\\[6pt]
{\small Professor da Rede Estadual de Ensino de São Paulo \textbar\ Pesquisador independente}\\[26pt]
{\small Versão 1.0 \ \textbar\ \ Outubro de 2026 \ \textbar\ \ Carga horária: 32 horas (8 encontros de 4 horas)}\\[10pt]
{\small Repositório: \url{https://github.com/ebiossanto/IA-RESEARCH-PESQUISA-ASSISTIDA}}
\vfill
{\small Bibliografia básica: WILSON (2001) \textbullet\ AMBROSE et al. (2010) \textbullet\ BOOTH et al. (2019) \textbullet\ UNESCO (2024)}
\end{titlepage}

\tableofcontents
\newpage

% ============================ PARTE A ============================
\section*{Parte A --- Carta de encaminhamento}
\addcontentsline{toc}{section}{Parte A --- Carta de encaminhamento}

{\small
\textbf{De:} Euzébio Soares dos Santos --- Professor da Rede Estadual de Ensino de São Paulo / Pesquisador independente\\
\textbf{Para:} Direção e Colegiado de Curso / Comissão de Ensino da instituição de destino\\
\textbf{Assunto:} Proposta de disciplina \emph{Pesquisa Assistida por Inteligência Artificial} (32 h) --- pedido de avaliação e de inclusão no currículo\\
\textbf{Local e data:} São Paulo, outubro de 2026
}\par\bigskip

Senhoras e Senhores Membros do Colegiado,

Encaminho, para avaliação, a proposta de disciplina \textbf{Pesquisa Assistida por
Inteligência Artificial}, de carga horária de 32 horas (8 encontros de 4 horas), destinada a
cursos de graduação, cursos técnicos, programas de extensão e à formação inicial de
pesquisadores.

O ponto de partida é simples e verificável: a IA generativa tornou fácil produzir textos,
códigos, imagens e análises preliminares, mas facilidade de produção não equivale a
compreensão, correção nem autoria. As instituições de ensino enfrentam hoje um dilema mal
resolvido --- proibir (o que é ineficaz e desloca o uso para fora da sala) ou permitir sem
critério (o que produz trabalhos plausíveis e não verificados). Esta disciplina propõe uma
terceira via: \textbf{ensinar metodologicamente o uso da IA dentro de um processo de
pesquisa}, com domínio conceitual, escolha consciente de ferramentas, verificação
independente e responsabilidade autoral.

A proposta organiza-se em quatro movimentos --- \textbf{DOMINAR, PREPARAR, SELECIONAR,
INVESTIGAR} --- e termina em dois produtos auditáveis: um ambiente computacional organizado e
um plano de investigação verificável. Não há memorização de marcas ou programas: o critério
de seleção de ferramentas é permanente (função, adequação, segurança, verificabilidade),
ainda que as ferramentas mudem.

A fundamentação teórica é composta por quatro obras de base, registradas na bibliografia
básica do documento:

\begin{itemize}
\item \textbf{UNESCO (2024)} orienta o uso crítico, humano e verificável da IA generativa na
educação e na pesquisa --- é o eixo da justificativa;
\item \textbf{AMBROSE et al. (2010)} sustenta a estrutura de aprendizagem baseada em pesquisa
(conhecimento prévio, prática com feedback, expectativas altas, aprendiz autônomo) --- é o
eixo da metodologia e da avaliação;
\item \textbf{WILSON (2001)} sustenta o domínio conceitual (definir, exemplificar, contraprovar,
relacionar) --- é o eixo das aulas 1 e 2;
\item \textbf{BOOTH, COLOMB e WILLIAMS (2019)} sustentam a passagem de interesse a problema, de
pergunta a hipótese e de hipótese a argumento com evidência --- é o eixo das aulas 7 e 8.
\end{itemize}

O documento encaminhado contém: a ementa acadêmica completa (Parte I), o documento
explicativo com justificativa, princípio pedagógico, protocolos de interação e auditoria
(Parte II) e a bibliografia básica com a articulação teórica (Parte III). O dossiê em PDF, a
síntese para comissão e o roteiro de apresentação acompanham esta carta.

Coloco-me à disposição para apresentação, ajuste de carga horária, adequação ao projeto
pedagógico do curso e para sanar dúvidas quanto à avaliação e aos requisitos de
infraestrutura.

Atenciosamente,

\textbf{Euzébio Soares dos Santos}\\
São Paulo --- Brasil

\newpage

% ============================ PARTE B ============================
\section*{Parte B --- Síntese para comissão de ensino}
\addcontentsline{toc}{section}{Parte B --- Síntese para comissão de ensino}

{\small
\textbf{Disciplina:} Pesquisa Assistida por Inteligência Artificial\\
\textbf{Carga horária:} 32 h (8 encontros de 4 h) \ \textbar\ \ \textbf{Natureza:} teórico-prática em laboratório\\
\textbf{Versão:} 1.0 --- Outubro de 2026 \ \textbar\ \ \textbf{Autor:} Euzébio Soares dos Santos
}\par\bigskip

\sub{1. Problema que a disciplina resolve}

{\small
\begin{tabular}{|p{0.30\textwidth}|p{0.31\textwidth}|p{0.31\textwidth}|}
\hline
\textbf{Situação atual} & \textbf{Consequência} & \textbf{Resposta desta disciplina} \\
\hline
IA generativa produz texto, código e imagem com facilidade & Produção sem compreensão; plausibilidade confundida com correção & Exigir domínio conceitual e histórico \emph{antes} de pedir respostas à IA \\
\hline
Uso proibido ou não orientado & Uso clandestino, trabalhos não verificáveis, sem critério de avaliação & Metodologia explícita de uso responsável, com registro e auditoria \\
\hline
Ferramenta escolhida por moda & Ambiente sobrecarregado, acúmulo sem finalidade & Seleção por função: entrada, saída, fontes, risco, privacidade, verificação \\
\hline
Plano de pesquisa improvisado & Pergunta vaga, hipótese não examinável, método sem evidência & Pergunta $\rightarrow$ hipótese $\rightarrow$ método $\rightarrow$ evidência $\rightarrow$ auditoria \\
\hline
\end{tabular}
}

\sub{2. Objetivo geral}

Capacitar o estudante a realizar uma investigação assistida por inteligência artificial,
partindo do domínio dos conceitos e da história de um tema, construindo um ambiente
computacional adequado, escolhendo as ferramentas e IAs apropriadas e elaborando uma
hipótese, um método e um plano verificável de pesquisa.

\sub{3. Os quatro movimentos da disciplina}

{\small
\textbf{1. DOMINAR} --- conceitos e história do problema (aulas 1--2) \ $\rightarrow$ \
\textbf{2. PREPARAR} --- ambiente, arquivos, versões e segurança (aula 3) \ $\rightarrow$ \
\textbf{3. SELECIONAR} --- ferramenta e IA adequadas a cada tarefa (aulas 4--5) \ $\rightarrow$ \
\textbf{4. INVESTIGAR} --- pergunta, hipótese, método, evidências e auditoria (aulas 6--8).
}

\sub{4. Grade resumida (8 encontros)}

{\small
\begin{tabular}{|p{0.06\textwidth}|p{0.44\textwidth}|p{0.44\textwidth}|}
\hline
\textbf{Aula} & \textbf{Tema} & \textbf{Produto entregue} \\
\hline
1 & Do tema aos conceitos & Ficha conceitual inicial \\
\hline
2 & História do problema e formação do método & Linha histórico-conceitual \\
\hline
3 & Construção do ambiente de pesquisa & Ambiente organizado e funcional \\
\hline
4 & Ferramentas para cada tarefa & Mapa pessoal de ferramentas \\
\hline
5 & Interação produtiva com IA & Registro comparativo de interações \\
\hline
6 & Auditoria, verificação e alucinações & Relatório de auditoria \\
\hline
7 & Pergunta, hipótese e método & Núcleo da investigação \\
\hline
8 & Plano integrado de investigação & Plano final de investigação \\
\hline
\end{tabular}
}

\sub{5. Avaliação e produto final}

{\small
\begin{tabular}{|p{0.38\textwidth}|p{0.12\textwidth}|p{0.44\textwidth}|}
\hline
\textbf{Componente} & \textbf{Peso} & \textbf{Evidência auditável} \\
\hline
Conceitos e história & 25\% & Ficha conceitual + linha histórico-conceitual \\
\hline
Ambiente e ferramentas & 20\% & Estrutura de pastas, diário de versões, mapa de ferramentas \\
\hline
Interação e auditoria da IA & 25\% & Registros de interação + relatório de auditoria com classificação \\
\hline
Plano de investigação & 30\% & Pergunta, hipótese principal e alternativa, método, critérios, cronograma \\
\hline
\end{tabular}
}

\textbf{Produto final:} dossiê compacto em quatro partes --- domínio do assunto; ambiente
preparado; estratégia de uso e auditoria de IA; pergunta, hipótese, método e plano de
investigação.

\sub{6. Requisitos mínimos}

\begin{itemize}
\item Laboratório de informática ou ambiente próprio dos estudantes (notebook com 4 GB de RAM).
\item Editor de texto, leitor de PDF, planilha e controle de versões (todos gratuitos).
\item Acesso a ao menos duas IAs distintas (uma textual e uma de código ou revisão crítica).
\item Nenhum software pago é requisito; nenhum dado pessoal de estudante sai do ambiente
institucional sem autorização.
\end{itemize}

\sub{7. O que a disciplina \emph{não} é}

Não é curso de ferramenta (nenhum produto comercial é objeto de avaliação); não é treinamento
de \emph{prompt} isolado (a instrução é uma etapa de um método maior); não é aula de
programação, estatística ou redação (entram como instrumentos do problema escolhido); não
substitui a orientação científica (encaminha o estudante a ela com trabalho verificável).

\caixa{\textbf{Justificativa institucional em uma frase:} a instituição não precisa escolher
entre proibir e permitir a IA: precisa ensinar, com critério e evidência, \emph{como} ela
pode ser usada --- e esta disciplina entrega esse ensinamento em 32 horas, com produto final
auditável.}

\newpage

% ============================ PARTE I ============================
\section*{Parte I --- Ementa acadêmica}
\addcontentsline{toc}{section}{Parte I --- Ementa acadêmica}

\subsection*{1. Identificação da disciplina}

{\small
\begin{tabular}{|p{0.30\textwidth}|p{0.64\textwidth}|}
\hline
\textbf{Nome} & Pesquisa Assistida por Inteligência Artificial \\
\hline
\textbf{Carga horária sugerida} & 32 horas \\
\hline
\textbf{Duração} & 8 encontros de 4 horas \\
\hline
\textbf{Natureza} & Teórico-prática, em formato de laboratório \\
\hline
\textbf{Público-alvo} & Estudantes de graduação, cursos técnicos, extensão e pesquisadores iniciantes \\
\hline
\textbf{Produto final} & Ambiente computacional organizado e plano de investigação assistida por IA \\
\hline
\end{tabular}
}

\subsection*{2. Ementa}

Introdução à pesquisa assistida por inteligência artificial. Estudo dos conceitos
fundamentais e da história de um problema. Distinção entre tema, conceito, problema,
pergunta, hipótese e método. Preparação de ambiente computacional para organização de
arquivos, fontes, dados, registros e versões. Configuração de agente computacional, como
OpenCode ou solução equivalente, com permissões adequadas. Seleção de aplicativos de apoio
para texto, referências, matemática, estatística, programação, imagem, diagramas e mapas
conceituais. Escolha da inteligência artificial conforme a natureza da tarefa. Técnicas de
interação produtiva, decomposição de problemas e comparação de respostas. Auditoria de
afirmações, referências, cálculos e resultados. Identificação de erros, contradições e
alucinações. Formulação de pergunta de pesquisa, hipótese, método e plano simples,
verificável e reprodutível de investigação.

\subsection*{3. Objetivo geral}

Capacitar o estudante a realizar uma investigação assistida por inteligência artificial,
partindo do domínio dos conceitos e da história de um tema, construindo um ambiente
computacional adequado, escolhendo as ferramentas e IAs apropriadas e elaborando uma
hipótese, um método e um plano verificável de pesquisa.

\subsection*{4. Objetivos específicos}

\begin{enumerate}
\item Identificar e compreender os conceitos fundamentais do tema escolhido.
\item Reconstruir de forma sintética a história do problema e de seus métodos.
\item Organizar no computador um ambiente de estudo e pesquisa seguro e reproduzível.
\item Selecionar programas de apoio de acordo com a necessidade real da tarefa.
\item Escolher uma ou mais IAs segundo a natureza do problema.
\item Interagir com a IA por meio de contexto, papéis, critérios e formatos claros.
\item Verificar afirmações, fontes, cálculos e resultados gerados ou auxiliados por IA.
\item Identificar alucinações, omissões, contradições e excesso de confiança.
\item Formular pergunta, hipótese, método e critérios de análise.
\item Produzir um plano de investigação documentado e executável.
\end{enumerate}

\subsection*{5. Competências}

{\small
\begin{tabular}{|p{0.30\textwidth}|p{0.64\textwidth}|}
\hline
\textbf{Competência} & \textbf{Descrição} \\
\hline
Domínio conceitual & Definir, relacionar e delimitar os conceitos essenciais do problema. \\
\hline
Compreensão histórico-metodológica & Reconhecer como o problema surgiu, mudou e foi investigado ao longo do tempo. \\
\hline
Preparação do ambiente & Organizar arquivos, versões, fontes, dados, ferramentas e registros de modo seguro. \\
\hline
Seleção de ferramentas & Escolher programas e IAs pela função necessária, evitando acúmulo sem finalidade. \\
\hline
Colaboração produtiva com IA & Usar IA para explicar, decompor, comparar, criticar, formalizar e revisar. \\
\hline
Auditoria e verificação & Distinguir plausibilidade de correção e verificar resultados de forma independente. \\
\hline
Planejamento investigativo & Manter coerência entre problema, pergunta, hipótese, método e evidência. \\
\hline
Responsabilidade autoral & Registrar o uso da IA e responder pelas decisões e conclusões apresentadas. \\
\hline
\end{tabular}
}

\subsection*{6. Resultados de aprendizagem}

Ao final da disciplina, o estudante será capaz de:

\begin{enumerate}
\item Explicar os conceitos centrais e a evolução histórica do problema escolhido.
\item Elaborar um glossário e um mapa de relações conceituais.
\item Configurar uma estrutura de projeto e um agente de IA com permissões controladas.
\item Selecionar aplicativos adequados para texto, matemática, estatística, imagem, mapas e dados.
\item Justificar a escolha de cada ferramenta e de cada IA.
\item Formular instruções claras, decompor tarefas e comparar respostas.
\item Auditar uma resposta de IA e classificar suas afirmações quanto à confiabilidade.
\item Identificar e corrigir pelo menos um erro, omissão ou alucinação.
\item Formular uma pergunta delimitada, uma hipótese examinável e uma hipótese alternativa.
\item Definir método, evidências, critérios de análise, limitações e cronograma.
\item Produzir um plano de investigação simples, documentado e reproduzível.
\end{enumerate}

\subsection*{7. Conteúdo programático em 8 aulas}

{\small
\begin{tabular}{|p{0.05\textwidth}|p{0.24\textwidth}|p{0.41\textwidth}|p{0.24\textwidth}|}
\hline
\textbf{Aula} & \textbf{Tema} & \textbf{Conteúdo central} & \textbf{Produto} \\
\hline
1 & Do tema aos conceitos & Tema, conceito, problema e pergunta; definições; exemplos; contraexemplos; glossário. & Ficha conceitual inicial. \\
\hline
2 & História do problema e formação do método & Origem dos conceitos; mudanças de definição; soluções anteriores; limitações e questões abertas. & Linha histórico-conceitual. \\
\hline
3 & Construção do ambiente de pesquisa & Pastas, arquivos, versões, diário, backup, privacidade, agente computacional e permissões. & Ambiente organizado e funcional. \\
\hline
4 & Ferramentas para cada tarefa & Programas de texto, referências, matemática, estatística, código, imagem, diagramas e mapas. & Mapa pessoal de ferramentas. \\
\hline
5 & Interação produtiva com IA & Contexto, papéis, decomposição, formato de saída, comparação de modelos e agentes especializados. & Registro comparativo de interações. \\
\hline
6 & Auditoria, verificação e alucinações & Checagem de fontes, cálculos, contradições, omissões, confiança e rastreabilidade. & Relatório de auditoria. \\
\hline
7 & Pergunta, hipótese e método & Delimitação; hipótese principal e alternativa; previsões; evidências; escolha do método. & Núcleo da investigação. \\
\hline
8 & Plano integrado de investigação & Ferramentas, IA, dados, critérios, riscos, cronograma, limitações e reprodutibilidade. & Plano final de investigação. \\
\hline
\end{tabular}
}

\subsection*{8. Metodologia}

A disciplina funcionará como laboratório. Em cada encontro, um conteúdo curto será
apresentado, demonstrado e imediatamente aplicado ao projeto do estudante. A ênfase estará no
processo de pesquisa, e não na memorização de marcas ou programas. As ferramentas poderão
mudar; o critério de seleção permanecerá baseado na função, na adequação, na segurança e na
verificabilidade.

Essa organização alinha-se aos princípios de ensino baseados em pesquisa: a aprendizagem
depende do que o estudante já sabe e de como organiza esse conhecimento, é sustentada por
metas claras, prática com feedback frequente e um clima que favoreça a tomada de
responsabilidade pelo próprio aprendizado (AMBROSE et al., 2010). Por isso cada aula termina
em \textbf{produto observável} e cada avaliação recai sobre trabalho produzido, não sobre
memorização.

\subsection*{9. Avaliação}

{\small
\begin{tabular}{|p{0.60\textwidth}|p{0.20\textwidth}|}
\hline
\textbf{Componente} & \textbf{Peso} \\
\hline
Conceitos e história & 25\% \\
\hline
Ambiente e ferramentas & 20\% \\
\hline
Interação e auditoria da IA & 25\% \\
\hline
Plano de investigação & 30\% \\
\hline
\end{tabular}
}

A avaliação é por portfólio: cada componente corresponde a um produto do dossiê final e pode
ser auditado de forma independente pela banca.

\subsection*{10. Produto final}

Um dossiê compacto composto por quatro partes: domínio do assunto; ambiente preparado;
estratégia de uso e auditoria de IA; pergunta, hipótese, método e plano de investigação.

\newpage

% ============================ PARTE II ============================
\section*{Parte II --- Documento explicativo}
\addcontentsline{toc}{section}{Parte II --- Documento explicativo}

\subsection*{1. Justificativa}

A inteligência artificial tornou fácil produzir textos, códigos, imagens e análises
preliminares. Isso, porém, não garante compreensão nem correção. A dificuldade educacional
não está apenas em permitir ou proibir a IA, mas em ensinar o estudante a utilizá-la dentro
de um processo de pesquisa que preserve domínio conceitual, escolha consciente de ferramentas,
verificação independente e responsabilidade humana.

O guia da UNESCO para a IA generativa na educação e na pesquisa (UNESCO, 2024) orienta no
mesmo sentido: a formação crítica e o acompanhamento humano devem vir antes da ferramenta ---
é preciso desenvolver no estudante a capacidade de usar a IA generativa com discernimento,
verificar o que ela produz e assumir a autoria do que apresenta. A disciplina é, portanto,
uma resposta institucional a essa demanda: em vez de regras de proibição, uma metodologia
explícita de uso responsável.

A disciplina propõe uma resposta prática: antes de solicitar respostas à IA, o estudante
compreende os conceitos e a história do problema; depois prepara seu ambiente de trabalho;
em seguida escolhe a ferramenta ou IA adequada; por fim formula e audita uma investigação.

{\small
\textbf{1. DOMINAR} --- conceitos, relações, história e limites do problema.\\
\textbf{2. PREPARAR} --- ambiente computacional, arquivos, programas e segurança.\\
\textbf{3. SELECIONAR} --- ferramenta e IA adequadas para cada tarefa.\\
\textbf{4. INVESTIGAR} --- pergunta, hipótese, método, evidências e auditoria.
}

\subsection*{2. O princípio pedagógico}

A IA pode ampliar a pesquisa, mas não substitui o domínio dos conceitos, a escolha do
método, a verificação dos resultados nem a responsabilidade pela conclusão.

O fluxo da disciplina é deliberadamente simples:

\fluxo

\subsection*{3. O que significa dominar os conceitos}

Dominar um conceito não significa decorar uma definição. Significa saber explicar o que ele
designa, como se relaciona com outros conceitos, onde se aplica, quais são seus exemplos e
contraexemplos, como surgiu e quais limites possui. É exatamente o programa de análise de
conceitos defendido por Wilson (2001): definir, exemplificar, contraprovar, relacionar e
organizar --- o conceito é dominado quando o estudante consegue operar com ele em situações
novas, não quando consegue repeti-lo.

Roteiro de verificação de domínio conceitual:

\begin{enumerate}
\item Qual é a definição adotada?
\item Existem definições concorrentes?
\item Quais exemplos e contraexemplos testam a definição?
\item Como o conceito mudou historicamente?
\item Que problema ele ajuda a formular?
\item O que ainda permanece incerto?
\end{enumerate}

\subsection*{4. O setup como laboratório pessoal}

O setup é a combinação organizada de arquivos, aplicativos, agentes e procedimentos que dá
suporte à investigação. Ele deve ser suficiente para o problema, mas não excessivo. O
estudante começa com uma estrutura mínima e acrescenta ferramentas somente quando surge uma
necessidade real.

\sub*{4.1 Estrutura sugerida}

{\small\begin{verbatim}
projeto-pesquisa/
|-- 01_conceitos/
|-- 02_historia/
|-- 03_fontes/
|-- 04_dados/
|-- 05_experimentos/
|-- 06_resultados/
|-- 07_imagens/
|-- 08_interacoes_ia/
|-- 09_auditoria/
`-- 10_relatorio/
\end{verbatim}}

\sub*{4.2 Categorias de ferramentas}

{\small
\begin{tabular}{|p{0.26\textwidth}|p{0.68\textwidth}|}
\hline
\textbf{Categoria} & \textbf{Finalidade} \\
\hline
Organização & Pastas, notas, controle de versões, backup e diário. \\
\hline
Texto e referências & Editor, Markdown, leitor de PDF e gerenciador bibliográfico. \\
\hline
Matemática & Calculadora simbólica, gráficos, equações e caderno computacional. \\
\hline
Estatística e dados & Planilha, ambiente estatístico e visualização. \\
\hline
Imagem e mapas & Editor de imagem, diagramas, fluxogramas, mapas mentais e conceituais. \\
\hline
Automação & Editor de código, terminal, controle de versões e agente computacional. \\
\hline
\end{tabular}
}

\subsection*{5. Como escolher a IA adequada}

Não é necessário utilizar a mesma IA para todas as tarefas. A escolha deve considerar o
tipo de entrada, o tipo de saída, a necessidade de fontes, o risco de erro, a privacidade
dos dados, a capacidade de cálculo ou código e a possibilidade de verificação.

\begin{itemize}
\item \textbf{Explicação conceitual:} priorizar clareza, comparação de definições e exemplos.
\item \textbf{Pesquisa documental:} priorizar rastreabilidade, fontes e acesso aos documentos.
\item \textbf{Matemática e estatística:} exigir passos, hipóteses e conferência em ferramenta independente.
\item \textbf{Programação:} utilizar ambiente com arquivos, testes, controle de versões e permissões.
\item \textbf{Imagem e diagramas:} definir objetivo visual e conferir fidelidade ao conteúdo.
\item \textbf{Revisão crítica:} utilizar uma IA ou agente separado para procurar falhas e alternativas.
\end{itemize}

\subsection*{6. Protocolo de interação produtiva}

Uma boa interação fornece à IA condições para compreender a tarefa e oferece ao estudante
condições para avaliar a resposta. O roteiro mínimo é:

\begin{enumerate}
\item Problema e contexto.
\item Conceitos já definidos.
\item Tarefa específica.
\item Restrições e ações proibidas.
\item Formato de saída esperado.
\item Fontes ou evidências necessárias.
\item Critérios de qualidade.
\item Procedimento de verificação.
\end{enumerate}

\subsection*{7. Auditoria e alucinações}

Uma alucinação é uma resposta apresentada com aparência de segurança, mas que contém
informação inventada, distorcida ou não sustentada. O estudante não deve auditar apenas o
texto final; deve separar as afirmações importantes e examinar cada uma.

{\small
\begin{tabular}{|p{0.42\textwidth}|p{0.52\textwidth}|}
\hline
\textbf{Pergunta de auditoria} & \textbf{Ação} \\
\hline
A afirmação é verificável? & Transformá-la em uma proposição clara. \\
\hline
A fonte existe? & Localizar o documento original. \\
\hline
A fonte sustenta a afirmação? & Ler o trecho e o contexto. \\
\hline
Há cálculo envolvido? & Refazer em ferramenta independente. \\
\hline
Há outra explicação? & Produzir hipótese alternativa. \\
\hline
Qual o grau de confiança? & Classificar e justificar. \\
\hline
\end{tabular}
}

\textbf{Classificação recomendada:} confirmado; parcialmente confirmado; contradito; não
verificável; incerto; deve ser reformulado.

\subsection*{8. Hipótese, método e plano de investigação}

A etapa final reúne todo o trabalho. A pergunta deve ser delimitada; a hipótese deve gerar
uma expectativa observável; e o método deve produzir evidências capazes de apoiar, limitar
ou contrariar a hipótese. É a passagem que Booth, Colomb e Williams (2019) descrevem como o
núcleo da pesquisa: transformar um interesse em problema, o problema em pergunta respondível
e a pergunta em uma resposta provisória --- a hipótese --- que só se sustenta quando
sustentada por evidência e argumento.

{\small
\begin{tabular}{|p{0.22\textwidth}|p{0.72\textwidth}|}
\hline
\textbf{Elemento} & \textbf{Pergunta correspondente} \\
\hline
Problema & Qual dificuldade ou lacuna será investigada? \\
\hline
Pergunta & O que exatamente se deseja descobrir? \\
\hline
Hipótese & Que resultado é esperado, em quais condições e por quê? \\
\hline
Método & Que procedimento poderá examinar a hipótese? \\
\hline
Evidência & Que observação apoiaria ou contrariaria a hipótese? \\
\hline
\end{tabular}
}

\subsection*{9. Modelo do plano final}

\begin{enumerate}
\item Título provisório.
\item Tema e delimitação.
\item Conceitos fundamentais.
\item Síntese histórica do problema.
\item Pergunta de pesquisa.
\item Hipótese principal e hipótese alternativa.
\item Método e procedimentos.
\item Dados, documentos ou fontes necessários.
\item Ferramentas e IAs selecionadas, com justificativa.
\item Critérios de auditoria e verificação.
\item Resultados possíveis e interpretação.
\item Riscos, limitações e cuidados éticos.
\item Cronograma de execução.
\item Registro de uso da IA.
\end{enumerate}

\subsection*{10. Conclusão}

A proposta não busca ensinar todas as áreas que podem participar de uma pesquisa. Ela
ensina o estudante a construir um caminho de investigação e a convocar os recursos
necessários no momento adequado. Matemática, estatística, programação, história, texto e
imagem entram como instrumentos do problema, não como blocos isolados.

O resultado esperado é um pesquisador iniciante capaz de \textbf{compreender antes de
automatizar, preparar antes de executar, escolher antes de utilizar e verificar antes de
concluir}.

\newpage

% ============================ PARTE III ============================
\section*{Parte III --- Bibliografia e articulação teórica}
\addcontentsline{toc}{section}{Parte III --- Bibliografia e articulação teórica}

\subsection*{Bibliografia básica}

\begin{enumerate}
\item AMBROSE, Susan A.; BRIDGES, Michael W.; DiPIETRO, Michele; LOVETT, Marsha C.; NORMAN,
Marie K. \textbf{How Learning Works: Seven Research-Based Principles for Smart Teaching}.
San Francisco: Jossey-Bass, 2010.

\item BOOTH, Wayne C.; COLOMB, Gregory G.; WILLIAMS, Joseph M. \textbf{A arte da pesquisa}.
3. ed. São Paulo: Martins Fontes, 2019.

\item UNESCO. \textbf{Guia para a IA generativa na educação e na pesquisa}. Paris: UNESCO, 2024.

\item WILSON, John. \textbf{Pensar com conceitos}. Tradução de Waldéa Barcellos. São Paulo:
Martins Fontes, 2001.
\end{enumerate}

\subsection*{Como a bibliografia sustenta a proposta}

{\small
\begin{tabular}{|p{0.22\textwidth}|p{0.47\textwidth}|p{0.25\textwidth}|}
\hline
\textbf{Obra} & \textbf{O que sustenta na disciplina} & \textbf{Onde aparece} \\
\hline
WILSON (2001) & Análise de conceitos: definir, exemplificar, contraprovar, relacionar e organizar mapas de conceitos. & Aulas 1--2; Parte II, seção 3 (domínio conceitual). \\
\hline
AMBROSE et al. (2010) & Sete princípios de ensino baseados em pesquisa: conhecimento prévio, organização do saber, motivação, prática com feedback, clima, expectativas altas e formação de aprendizes autônomos. & Parte I, seções 8--9 (metodologia e avaliação); sequência das 8 aulas. \\
\hline
UNESCO (2024) & Uso crítico e humano da IA generativa em educação e pesquisa; verificação, ética e autoria. & Parte II, seção 1 (justificativa); Parte I, seção 2 e Parte II, seções 5--7. \\
\hline
BOOTH et al. (2019) & Do interesse ao problema, da pergunta à hipótese, do argumento à evidência e à apresentação dos resultados. & Aulas 7--8; Parte II, seções 8--9 (pergunta, hipótese, método e plano). \\
\hline
\end{tabular}
}

\subsection*{Como citar este trabalho}

{\small
SOARES, Euzébio. \textbf{Pesquisa Assistida por Inteligência Artificial: ementa acadêmica e
documento explicativo}. Versão 1.0. São Paulo, 2026. Disponível em:
\url{https://github.com/ebiossanto/IA-RESEARCH-PESQUISA-ASSISTIDA}.
}

\end{document}
